%% body.tex -- main content of the paper 
%% Use \rev{...} or \begin{revision}...\end{revision} for changes made
%% during peer review; switch colour on/off in main.tex.

%% =====================================================================
\section{Introduction}\label{sec:intro}
%% =====================================================================
Introduce the application domain and the decision or modeling problem. Explain why existing methods are insufficient~\citep{boyd2004convex,nocedal2006numerical}. State the research question.

\rev{This sentence shows how a change requested by a reviewer appears. Set \texttt{\textbackslash highlightrevisionsfalse} in \texttt{main.tex} to typeset it in black.}

Our contributions are:
\begin{itemize}
  \item \textbf{Modeling.} A new formulation of \ldots
  \item \textbf{Theory.} We prove that \ldots\ (\cref{thm:main}).
  \item \textbf{Algorithm.} An efficient procedure with complexity \ldots\ (\cref{alg:method}).
  \item \textbf{Experiments.} Numerical evidence on \ldots\ (\cref{sec:results}).
\end{itemize}
The remainder of the paper is organized as follows. \Cref{sec:related} reviews related work, \cref{sec:problem} formulates the problem, \cref{sec:method} presents the method, \cref{sec:results} reports numerical results, and \cref{sec:conclusion} concludes.

%% =====================================================================
\section{Related Work}\label{sec:related}
%% =====================================================================
\paragraph{Topic A.} Summarize the first line of work~\citep{tibshirani1996regression}.

\paragraph{Topic B.} Summarize the second line of work~\citep{rasmussen2006gaussian,benner2015survey}.

\paragraph{Positioning.} Explain how this paper differs from the works above.

%% =====================================================================
\section{Problem Formulation}\label{sec:problem}
%% =====================================================================
\paragraph{Notation.} Vectors are bold lowercase ($\bm{x}\in\R^n$), matrices bold uppercase ($\bm{A}\in\R^{m\times n}$), $\norm{\cdot}$ denotes the Euclidean norm, $\inner{\bm{x}}{\bm{y}}=\bm{x}\T\bm{y}$, and $[n]\coloneqq\{1,\dots,n\}$.

We consider the problem
\begin{mini}|s|
  {\bm{x}\in\R^n}{f(\bm{x}) \coloneqq \E_{\xi}\big[F(\bm{x},\xi)\big]}{\label{eq:problem}}{}
  \addConstraint{g_i(\bm{x})}{\le 0,}{\quad i\in[m]}
  \addConstraint{\bm{A}\bm{x}}{=\bm{b},}{}
\end{mini}
where $F:\R^n\times\Xi\to\R$ is \ldots\ and $\xi$ is a random vector with distribution $\Prob$.

\begin{assumption}\label{ass:smooth}
The function $f$ is convex and $L$-smooth, i.e., $\norm{\nabla f(\bm{x})-\nabla f(\bm{y})}\le L\norm{\bm{x}-\bm{y}}$ for all $\bm{x},\bm{y}$.
\end{assumption}

%% =====================================================================
\section{Methodology}\label{sec:method}
%% =====================================================================
Describe the proposed approach. A displayed equation with a label:
\begin{equation}\label{eq:update}
  \bm{x}_{k+1} = \proj_{\mathcal{X}}\big(\bm{x}_k - \eta_k \nabla f(\bm{x}_k)\big).
\end{equation}

\begin{algorithm}[t]
  \caption{Proposed method}\label{alg:method}
  \begin{algorithmic}[1]
    \Require initial point $\bm{x}_0$, step sizes $\{\eta_k\}$, tolerance $\varepsilon$
    \For{$k = 0,1,2,\dots$}
      \State compute $\nabla f(\bm{x}_k)$
      \State update $\bm{x}_{k+1}$ by \eqref{eq:update}
      \If{$\norm{\bm{x}_{k+1}-\bm{x}_k}\le\varepsilon$} \State \textbf{break} \EndIf
    \EndFor
    \Ensure $\bm{x}_{k+1}$
  \end{algorithmic}
\end{algorithm}

\subsection{Theoretical Analysis}\label{sec:theory}
\begin{theorem}\label{thm:main}
Under \cref{ass:smooth}, with $\eta_k = 1/L$, the iterates of \cref{alg:method} satisfy
\[
  f(\bm{x}_k) - f^\star \le \frac{L\norm{\bm{x}_0-\bm{x}^\star}^2}{2k}.
\]
\end{theorem}
\begin{proof}
Sketch the argument or refer to the appendix.
\end{proof}

\begin{remark}
Comment on the assumptions, tightness of the bound, or extensions.
\end{remark}

%% =====================================================================
\section{Numerical Experiments}\label{sec:results}
%% =====================================================================
\subsection{Experimental Setup}\label{sec:setup}
Describe datasets or test instances, baselines, hyperparameters, hardware, and the number of repetitions. Report tolerances with \texttt{siunitx}, e.g., $\varepsilon=\num{1e-6}$ and a time limit of \SI{3600}{\second}.

\subsection{Results}\label{sec:main-results}
\Cref{tab:results} summarizes the main results and \cref{fig:placeholder} shows \ldots

\input{table}
\input{fig}

\begin{revision}
A longer revised passage, including equations, can be wrapped in the \texttt{revision} environment:
\begin{equation}
  \mathrm{Gap}_k = \frac{f(\bm{x}_k)-f^\star}{\abs{f^\star}}.
\end{equation}
\end{revision}

\subsection{Sensitivity Analysis}\label{sec:sensitivity}
Discuss robustness to parameters, problem size, or noise.

%% =====================================================================
\section{Discussion}\label{sec:discussion}
%% =====================================================================
Interpret the results, state limitations, and outline when the method should or should not be used.

%% =====================================================================
\section{Conclusion}\label{sec:conclusion}
%% =====================================================================
Summarize the contributions and findings, and give directions for future work.

\section*{Acknowledgments}
Funding sources and acknowledgments (remove for anonymous submissions).
